%%%%%%%%%%%%%%%%%%%%%%%%%%%%%%%%%%%%%%%%%%%%%%%%%%%%%%%%%%%%%%%%%%%%%%%%%%%%%%%%
%% 07_Documentacao_Treinamento.tex
%% Guias para os Usuários: Operação e Administração do Sistema
%%%%%%%%%%%%%%%%%%%%%%%%%%%%%%%%%%%%%%%%%%%%%%%%%%%%%%%%%%%%%%%%%%%%%%%%%%%%%%%%

\chapter{Guias para os Usuários}
\label{cap:guias_usuarios}

\section{Objetivo dos Guias}

Este capítulo apresenta as instruções de operação da plataforma NextGIS Web e dos serviços de rede. O conteúdo é organizado para atender aos administradores técnicos da infraestrutura e aos pesquisadores do grupo de Modelagem Hidrológica.

\section{Guia do Administrador: Implantação e Manutenção}

\subsection{Inicialização dos Serviços}

A implantação do ambiente nos servidores locais ou na instância de treinamento em nuvem requer a instalação prévia do Docker e do Docker Compose. A inicialização e verificação dos serviços executam os seguintes comandos:

\begin{enumerate}
    \item Acessar o diretório de implantação da aplicação:
\begin{lstlisting}[language=bash]
cd deploy/
\end{lstlisting}
    \item Inicializar os contêineres em segundo plano:
\begin{lstlisting}[language=bash]
docker compose up -d
\end{lstlisting}
    \item Verificar o estado de execução dos serviços:
\begin{lstlisting}[language=bash]
docker compose ps
\end{lstlisting}
    Os contêineres \texttt{ngw\_postgis}, \texttt{ngw\_redis}, \texttt{ngw\_core} e \texttt{ngw\_proxy} devem apresentar o estado \texttt{Up}.
\end{enumerate}

\section{Guia do Pesquisador: Uso da Plataforma Web GIS}

O acesso dos pesquisadores e estudantes à plataforma ocorre diretamente pelo navegador web de forma simplificada, sem necessidade de interação com o terminal de comandos.

\subsection{Acesso e Login}
\begin{enumerate}
    \item Abrir o navegador e acessar o endereço institucional do laboratório (\url{http://nextgis.labhidro.ufs.br} ou o endereço IP público fornecido pelo administrador durante os treinamentos);
    \item Clicar no botão \textbf{Entrar} no canto superior direito;
    \item Inserir o nome de usuário e a senha fornecidos pela equipe de suporte do laboratório.
\end{enumerate}

\subsection{Importação de Camadas Vetoriais e Matriciais}
\begin{enumerate}
    \item No painel lateral, selecionar a pasta da bacia hidrográfica correspondente ao estudo;
    \item Clicar em \textbf{Criar Recurso} e selecionar o tipo de dado: \textbf{Camada Vetorial} para polígonos ou linhas de drenagem, ou \textbf{Camada Raster} para Modelos Digitais de Elevação;
    \item Fazer o upload do arquivo compactado em formato \texttt{.zip} contendo os arquivos do Shapefile (\texttt{.shp}, \texttt{.shx}, \texttt{.dbf} e \texttt{.prj}) ou arquivo \texttt{.geojson};
    \item Confirmar o Sistema de Referência Espacial detectado (SIRGAS 2000 / UTM 24S - EPSG:31984) e salvar o recurso.
\end{enumerate}

\subsection{Criação de Mapas Interativos}
\begin{enumerate}
    \item No menu de recursos, clicar em \textbf{Criar Recurso} $\rightarrow$ \textbf{Mapa Web};
    \item Marcar as camadas vetoriais de hidrografia e saneamento a serem exibidas sobre o mapa base (OpenStreetMap ou relevo);
    \item Definir a escala inicial de visualização e salvar o mapa;
    \item O link gerado pode ser compartilhado com os demais integrantes da equipe para consulta rápida.
\end{enumerate}

\subsection{Conexão com o QGIS Desktop}
Para análises hidrológicas avançadas em estações de trabalho locais, os dados podem ser consumidos diretamente no software QGIS:
\begin{enumerate}
    \item No NextGIS Web, abrir as propriedades da camada desejada e copiar o link do serviço \textbf{WMS} (\textit{Web Map Service});
    \item No QGIS, clicar no menu \textbf{Camada} $\rightarrow$ \textbf{Adicionar Camada} $\rightarrow$ \textbf{Adicionar Camada WMS/WMTS};
    \item Criar uma nova conexão colando o link copiado;
    \item Clicar em \textbf{Conectar} e adicionar as feições ao projeto local.
\end{enumerate}
